%% kapa: exam=1 round=29 year=2018 period=1 form=A subject=civil_law subject_ko=민법 seq=26 number=26
%% source: Q-Net 감정평가사 기출문제 (cst00309) artlSeq=5202808 file=extracted/1008/5202808/2018년 제29회 감정평가사 1차시험 기출문제/2018년 제29회 감정평가사 1차시험 기출문제/1교시 A형.pdf
%% conversion: kordoc 4.24.51 -> tools/gen1.py (auto); review=proofread-sonnet; flags=none
\begin{question}{26}
저당권에 관한 설명으로 옳지 않은 것은? (다툼이 있으면 판례에 따름)
\begin{choices}
  \item 저당물의 소유권을 취득한 제3자는 경매인이 될 수 없다.
  \item 토지를 목적으로 저당권을 설정한 후 그 설정자가 그 토지에 건물을 축조하고 소유한 경우, 저당권자는 토지와 함께 그 건물에 대하여도 경매를 청구할 수 있다.
  \item 저당부동산에 대하여 저당권에 기한 압류가 있으면, 압류 이후의 저당권설정자의 저당부동산에 관한 차임채권에도 저당권의 효력이 미친다.
  \item 저당부동산에 대하여 지상권을 취득한 제3자는 저당권자에게 그 부동산으로 담보된 채권을 변제하고 저당권의 소멸을 청구할 수 있다.
  \item 저당권설정자의 책임있는 사유로 인하여 저당물의 가액이 현저히 감소된 때에는 저당권자는 저당권설정자에 대하여 그 원상회복 또는 상당한 담보제공을 청구할 수 있다.
\end{choices}
\end{question}
\answer{1}
